Three independent tabular Q-learners bid repeatedly into a stylised uniform-price electricity auction with three symmetric 60~MW firms, marginal cost \eur{10}/MWh and inelastic demand of 100~MW. After 5M rounds every seed has converged and the mean collusion index $\Dl$ is 0.621 (95\% CI 0.612--0.630), and, with firm 1 forced to undercut for one round, rivals' bids fall in 91\% of seeds but return to their old level within 24 rounds in only 21\%; the per-seed verdict changes with the deviator's identity in 61 of 100 seeds. The standard indicators of tacit collusion nevertheless fail to attribute the price to it. Removing the future ($\gamma = 0$) lowers $\Dl$ to 0.370, above the pre-registered collapse threshold of 0.25; removing memory raises it to 0.943. An unplayed profitable deviation exists in 100\% of seeds, but also in 100\% of memory-0 control seeds. Four further pre-registered experiments, two of whose predictions failed, are consistent with a verbal mechanism: the memory-0 price is the small-noise limit of an $\varepsilon$-greedy Edgeworth-type cycle. It survives ten times as many exploratory rounds, a restart of exploration from $\varepsilon = 0.5$, and a change to pay-as-bid. A fifth, a Q-table test, the share of temptations where a single Bellman backup would already flip the decision, is 100\% in every memory-0 seed and at most 80\% in any full-learner seed (mean 31\%). It separates a decision that rests on a learned continuation from one that rests on a stale estimate; the behavioural indicators do not.
